\section{OpenMP--MPI Implementation}

The CPU-parallel backend follows the natural coarse-grained decomposition of ACO.  OpenMP threads build independent ant solutions inside each rank, while MPI combines information across ranks.  This mapping keeps the sequential decision chain of an individual ant local and exposes parallelism across ants.

\subsection{Shared-memory organization}

The current implementation creates one persistent OpenMP region around the epoch loop.  Thread-private route-construction workspaces avoid sharing the visited set and random-number state.  Candidate indices and heuristic scores are stored contiguously, while the distance and pheromone matrices use single-precision aligned storage within the parallel solver.  Per-epoch work consists of four main phases: candidate-score refresh, independent ant construction, selection of the best solution, and pheromone evaporation/deposit.

The ant loop uses a runtime-selectable OpenMP schedule.  Pheromone evaporation is partitioned statically over the dense matrix; deposits use atomic additions because different routes can touch the same edge.  This removes a serial deposit phase at the price of possible atomic contention.  The fixed-size strong-scaling campaign in Section~\ref{sec:performance} is designed to measure the net effect rather than attributing time to any one phase without profiling.

\subsection{Distributed-memory synchronization}

Each MPI rank owns a replica of the pheromone matrix and receives a disjoint portion of the total ant population.  Modified edges are represented as sparse \emph{delta} records containing a matrix index and an increment.  The implementation starts a nonblocking \texttt{MPI\_Iallgatherv}; deltas from the previous epoch are applied before the next local update.  This avoids transmitting the complete dense matrix when only route edges have changed.

The sparse exchange does not remove all distributed overhead.  Every rank still initializes and evaporates a dense matrix, performs collective coordination, and stores its local replica.  Consequently, increasing the number of ranks may expose replicated work even when the communicated payload is sparse.  The observed multi-node trend is discussed conservatively in Section~\ref{sec:performance}; no network or phase-level profile is available in the current result set.

\subsection{Alternatives that were rejected}

Earlier prototypes in \texttt{experiments/openmp\_v2/} explored row-wise dense collectives and collaborative \emph{intra-ant} execution.  The former issued many collectives per epoch.  The latter assigned several CPU threads to one ant, using barriers, spinning counters, or condition variables to coordinate each construction step.  These approaches increased the coordination granularity and were abandoned in favor of one independent ant per worker.

The repository does not contain compatible raw timing files for those prototypes, so their previously quoted speedups are not reused in the final quantitative evaluation.  The implementation history nevertheless motivates the final design choice: independent ants need synchronization only at phase boundaries, whereas collaborative ants synchronize along the sequential route-construction path.  SIMD within one core remains a possible alternative because it could accelerate candidate scans without introducing inter-core coordination.

\subsection{Correctness boundary}

The current validation reconstructs every saved route and checks customer coverage and capacity.  All available OpenMP--MPI routes are feasible.  For several multi-rank runs, however, the scalar \texttt{best\_cost} does not match the cost of the route printed by rank zero.  This does not change the fixed-epoch wall time, which is retained for scaling, but it prevents using those cost values as multi-rank quality evidence.  Resolving the ownership and transfer of the globally best route is therefore a necessary correctness improvement beyond the timing study.
